\label{chap:p3-83-46-godel_turing}

\section{Effective decidability of each finite window of projective multiplicities}

\begin{proposition}
At a fixed depth \(N\), if the domains are finite and the atomic
observables are computable, every formula in the finite language
\(L_N\) is semantically decidable.
\end{proposition}

\begin{proof}
Atomic formulas are computed; connectives are decided by truth
tables; each quantifier ranges over a finite domain. The conclusion follows by
induction on syntactic complexity.
\end{proof}

The decidability of each window does not necessarily provide a single
algorithm that decides a property of the complete tower.

\section{Effective tower associated with the computation of projective multiplicities}

Let \(T_e\) be the Turing machine with index \(e\). For each \(n\), define
\[
 X_n^e
 =
 \{b_n\}
 \cup
 \{c_n:T_e\text{ has not halted in its first }n\text{ steps}\}.
\]
Links preserve existing labels:
\[
  p_{n+1,n}(b_{n+1})=b_n,
  \qquad
  p_{n+1,n}(c_{n+1})=c_n.
\]
The visible tower has to single symbol \(v_n\), and
\[
  q_n(b_n)=q_n(c_n)=v_n.
\]

\begin{theorem}[Global multiplicity and stopping]
\label{p3-83-c46-s1:thm:halt-fibre}
\[
 \#q_\infty^{-1}(\boldsymbol v)
 =
 \begin{cases}
 1,&T_e\text{ stops},\\
 2,&T_e\text{ never halts}.
 \end{cases}
\]
Therefore, no total algorithm decides uniqueness of the global fiber for
all effective towers in this family.
\end{theorem}

\begin{proof}
The branch \(\boldsymbol b=(b_n)\) always exists. If \(T_e\) halts at
step \(h\), the states \(c_n\) disappear from \(h\) onward and do not form a
global section. If it never halts, \(\boldsymbol c=(c_n)\) is a second
section. A decider for uniqueness would decide \(\HALT\), contradicting
\citep{turing1936}.
\end{proof}

We define
\[
  U_e:\#q_\infty^{-1}(\boldsymbol v)=1.
\]
Then \(U_e\iff T_e\downarrow\).

\begin{theorem}[Uniform incompleteness]
\label{p3-83-c46-s1:thm:uniform-incompleteness}
Let \(\mathsf T\) be a recursively axiomatizable theory that formalizes the
towers \(X^e\) and is semantically correct for \(U_e\) and
\(\neg U_e\). Then \(\mathsf T\) does not decide every \(U_e\).
\end{theorem}

\begin{proof}
If each \(U_e\) were to decide, for an index \(e\) one would enumerate in parallel
the proofs of \(U_e\) and \(\neg U_e\) until finding one. By correctness,
the first would appear exactly when the machine halts and the second
when it does not. A decider of \(\HALT\) would be obtained.
\end{proof}

\begin{figure}[H]
\centering
\begin{tikzpicture}[>=Latex,node distance=11mm and 15mm,
  box/.style={rounded corners,draw=hmtblue,fill=hmtlightblue,
    minimum width=39mm,minimum height=13mm,align=center},
  warn/.style={rounded corners,draw=hmtred,fill=hmtlightred,
    minimum width=39mm,minimum height=13mm,align=center}]
  \node[box] (sim) {finite simulation of \(T_e\)};
  \node[box,below=of sim] (levels) {\(X_n^e\) computable\\for every \(n\)};
  \node[box,below=of levels] (fibre)
    {\(\#q_\infty^{-1}(v)=1\iff T_e\downarrow\)};
  \node[warn,below left=of fibre] (undec) {no uniform decider};
  \node[warn,below right=of fibre] (inc)
    {sound effective theory\\incomplete for \(U_e\)};
  \draw[->,thick,hmtblue] (sim)--(levels);
  \draw[->,thick,hmtblue] (levels)--(fibre);
  \draw[->,thick,hmtred] (fibre)--(undec);
  \draw[->,thick,hmtred] (fibre)--(inc);
\end{tikzpicture}
\caption{From the finite decision to uniform incompleteness.}
\label{p3-83-c46-s1:fig:godel-turing}
\end{figure}

\section{Finite completeness of the HMT levels and exact domain of the Gödel--Turing obstruction}

The academic name used is:
\[
\boxed{\text{HMT uniform incompleteness theorem for projective multiplicities}.}
\]
The abbreviated designation “third Gödel result in HMT” is also used. The proof
is Gödel--Turing: it does not reproduce the self-referential diagonal of
\cite{goedel1931}, but instead reduces HALT to the multiplicity of a fibre.

\begin{remark}[Conditions of validity]
The theorem does not prove the inconsistency of ZFC. Nor does it preclude a
direct proof that a particular tower has a unique section. It precludes a
universal uniqueness procedure for all effective towers.
\end{remark}

% REMATE_LOCAL_GODEL_CONCLUSION
\Needspace{25\baselineskip}
\section{Exact relationship with the \(9\) phases of the Topological Phase Kernel}

The \(9\) phases of the TPK provide a concrete algorithm that, with its declared readers,
partitions the refinement block into \(729\) subcylinders and selects one of them
at each finite stage.
The passage to the infinite section uses a global theorem of nonemptiness,
compatibility, and unitarity.

The Gödel--Turing result says:
\[
\text{there is no uniqueness decider for every effective tower}.
\]
It does not state:
\[
\text{no concrete tower can have a uniform proof of its own}.
\]
And a particular proof does not provide the universally forbidden decisor.


\begin{remark}[Structural Conclusion]
Finite decidability, existence of the limit, uniqueness of an instance, and uniform
decidability of all instances are four levels. Holographic closure
requires that they be declared, rather than automatically transferring a property from one
level to another.
\end{remark}

\begin{remark}[Separation of results and scope]
The reducing tower, the equivalence \(U_e\iff T_e\downarrow\), and the metamathematical
theorem are \emph{results established in the preceding sections}. The central position within
the closure of infinity is an expository organization that alters neither the logical content nor the theorem's domain.
\end{remark}
